\documentclass{article}
\usepackage[utf8]{inputenc}
\usepackage{hyperref}
\usepackage{fullpage}


\newcommand{\eps}{\epsilon}
\newcommand{\light}{\textsf{light}}
\newcommand{\heavy}{\textsf{heavy}}
\renewcommand{\ll}{\textsf{$\ell-\ell$}}
\newcommand{\hh}{\textsf{$h-h$}}
\newcommand{\lh}{\textsf{$$\ell-h$}}
\newcommand{\poly}{\textrm{poly}}

\begin{document}



\subsection*{Dealing with 4-cycles on very-high degree vertices}
Let $\theta = t n^{1/2}$ for $t \geq 2$. We denote by $H^*$ the set of vertices with degree at least $\theta$, which we refer to as \emph{very-high} degree vertices. We also have the set $L$ of low-degree vertices, whose degree is at most $\alpha$ (or $\alpha/\epsilon$ or so) and the set $H$ of high-degree vertices (the remaining vertices).

Assume we are in the case that we have $\Omega(\eps m)$ edge-disjoint $C_4$s that each reside on two vertices in $L$ and two in $H^*$, and furthermore, each vertex $v$ that participates in these cycles has $\Omega(\eps d(v))$ edges that participate in them. We'll make one simplifying assumption (which should be possible to remove at a cost of $\poly(1/\epsilon)$) that we have a bipartite graph with the vertices of $L$ on one side and the vertices of $H^*$ on the other (and no other vertices an edges).

We will make another assumption that will need dealing with based on how the arboricity comes into play in the graph, by which all vertices in $L$ have degree $\alpha$ (or so). For some things it will be sufficient to use the fact that all vertices in $H^*$ have degree at least $\theta$, but there is one place that for now I'll need degree exactly $\theta$, so this will be another thing we'll need to deal with. What we get from the fact that all vertices in $H^*$ have degree at least $\theta$, is that
\begin{equation}\label{eq:Hst-lb}
|H^*| \leq \frac{m}{\theta} = \frac{m}{tn^{1/2}} \leq \frac{1}{t}\cdot n^{1/2}\alpha\;,
\end{equation}
so that
\begin{equation}\label{eq:Hst-ratio}
\frac{|H^*|}{\theta\alpha} \leq \frac{m/\theta}{\theta \alpha} = \frac{m}{t^2 n \alpha} = \frac{1}{t^2}\cdot \frac{m}{n\alpha} \leq \frac{1}{t^2}\;.
\end{equation}


We want to show that for any vertex $v$ in $H^*$, if we take something like $|H^*|^{1/2}$ random walks of length-2 from $v$, then we'll observe a $C_4$ with high constant probability, i.e., we'll get one walk $v-u-v'$ and one walks $v-u'-v'$ (where $u,u'$ are neighbors of $v$ in $L$ and $v' \in H^*$).

Let $u_1,\dots,u_\theta$ be the neighbors of $v$, where here I am assuming that it has $\theta$ neighbors and not more. Since we'll need to deal with vertices in $H^*$ that have varying degrees, we'll probably try to exploit the fact that we are selecting these vertices with probability proportional to their degree.

Let $v_1,\dots,v_b$ be the vertices in $H^*$ with which $v$ participates in $C_4$'s. For each $j\in [b]$, let $S_j$ be the subset of neighbors of $v$ that also neighbor $v_j$. Since $v$ and $v_j$ participate in at least one common $C_4$, this means that $|S_j| > 1$ for every $j \in [b]$. Let $R$ be the $r$ vertices in $H^*$ such that there is a single length-2 path from $v$ to each of them.
We have that
\begin{equation}\label{eq:b-r-Sj}
b+r \leq  H^* \;\;\mbox{ and}\;\; \sum_{j=1}^b |S_j| + r = \theta \alpha\;.
\end{equation}


In order to hit a $C_4$ in two distinct random walks, the first walk should go from $v$ to some neighbor $u_i$ of $v$, and from $u_i$ to some $v_j$ ($j \in [b])$, and then the second walk should go from $v$ to some other $u_k \in S_j$ (for the same $j$), and then from $u_k$ to $v_j$. The probability of this event is
\begin{equation}
\frac{1}{\theta}\sum_{i=1}^{\theta} \sum_{j\in [b]: u_i \in S_j} \frac{1}{\alpha} \cdot \frac{|S_j|-1}{\theta}\cdot \frac{1}{\alpha} \;=\; \frac{1}{(\theta\alpha)^2} \cdot \sum_{j=1}^b |S_j|\cdot (|S_j|-1)
\; \geq \frac{1}{(\theta\alpha)^2} \cdot \sum_{j=1}^b \frac{1}{2}|S_j|^2\;,
\label{eq:col-prob}
\end{equation}
(recall that $|S_j|>1$ for each $j \in [b]$).

What do we know about the sum of the $S_j$'s? Since $r < |H^*|$, we have that
\begin{equation}\label{eq:sym-of-Sj}
\sum_{j=1}^b |S_j| > \theta \alpha - |H^*| \geq (\theta\alpha)(1-1/t^2)\;.
\end{equation}
Since conditioned on the sum of the $|S_j|$s, the expression in Equation~(\ref{eq:col-prob}) is minimized when they are all of equal size, we get that the collision probability is at least
\begin{equation}
\frac{1}{(\theta\alpha)^2} \cdot \frac{1}{2}\cdot b\cdot \frac{(\theta\alpha)^2}{b^2}\cdot \frac{(t^2-1)^2}{t^4}
\;\geq\; \frac{1}{c|H^*|}\;,
\label{eq:col-prob2}
\end{equation}
for some small constant $c$.

\subsection*{trying to remove some assumptions}

Slightly changing notation, let $\theta^* = tn^{1/2}$ be the minimum degree of vertices in $H^*$ (for $t \geq 2$), and let $\theta \geq \theta^*$ be the degree of some vertex $v$ in $H^*$. Let $u_1,\dots,u_\theta$ be the neighbors of $v$ in $L$ and $\alpha_i$ be the degree of  $u_i$, so that $\alpha_i \leq \alpha$ (or $\alpha/\eps$). The subsets $S_j$, $j \in [b]$ are defined as above, and so is $R$.
The number of (distinct) length-2 walks from $v$ is $\sum_{i=1}^\theta \alpha_i$, so that 
\begin{equation}\label{eq:num-of-walks-gen}
\sum_{j=1}^b |S_j| + r = \sum_{i=1}^\theta \alpha_i\;.
\end{equation}
We also know that if we sum over the degrees of all vertices in $L$, then we get $m$ (under our simplifying assumption of a bipartite graph).

The probability of a collision of two length-2 walks when starting from $v$ is
%
\begin{equation}
\frac{1}{\theta}\sum_{i=1}^{\theta}  \sum_{j\in [b]: u_i \in S_j} \frac{1}{\alpha_i} \cdot 
\sum_{k\neq i: u_k \in S_j} \frac{|S_j|-1}{\theta}\cdot \frac{1}{\alpha_k} \;=\; \frac{1}{\theta^2} 
\cdot \sum_{j=1}^b \sum_{i\neq k: u_i,u_k\in S_j} \frac{1}{\alpha_i\cdot\alpha_k}
%\; \geq \frac{1}{(\theta\alpha)^2} \cdot \sum_{j=1}^b \frac{1}{2}|S_j|^2\;,
\label{eq:col-prob-gen}
\end{equation}

I am not sure how to continue from here. Actually, maybe we can't give a lower bound on the collision probability starting from every $v \in H^*$, because it is possible that for some such $v$ many of its neighbors in $L$ have degree much lower than $\alpha$, so that the ratio between the number of length-2 walks from $v$ and $H^*$ is not large as we need it to be. As you said, if $m$ is indeed roughly $n\alpha$, then maybe we can just ignore such vertices $v$ (?)
In general we might need to discuss what happens when we select $v$ at random (with probability proportional to its degree), and this may come up also with 4-cycles of the form $L-H-L-H$.




\end{document}
